\documentclass[11pt]{article}

\usepackage[a4paper,margin=25mm]{geometry}
\usepackage{fontspec}
\usepackage{amsmath,amssymb,mathtools}
\usepackage{booktabs,array,tabularx}
\usepackage{enumitem}
\usepackage{xcolor}
\usepackage{hyperref}
\usepackage{fancyhdr}
\usepackage{tcolorbox}
\usepackage{listings}
\usepackage{microtype}

\definecolor{navy}{HTML}{17365D}
\definecolor{blue}{HTML}{2D5ABB}
\definecolor{lightblue}{HTML}{EEF4FB}
\definecolor{warmgray}{HTML}{F5F4F1}
\definecolor{darkgray}{HTML}{444444}

\hypersetup{colorlinks=true,linkcolor=navy,urlcolor=blue,citecolor=navy}
\setlength{\parindent}{0pt}
\setlength{\parskip}{0.55em}
\setlength{\headheight}{14pt}
\setlength{\emergencystretch}{2em}
\setlist{itemsep=0.25em,topsep=0.35em}
\pagestyle{fancy}
\fancyhf{}
\lhead{Computation and Abstraction}
\rhead{Lecture Notes}
\cfoot{\thepage}
\renewcommand{\headrulewidth}{0.4pt}

\newtcolorbox{keyidea}[1][]{colback=lightblue,colframe=blue,
  boxrule=0.7pt,arc=1.5mm,title={#1},fonttitle=\bfseries}
\newtcolorbox{definition}[1][]{colback=warmgray,colframe=darkgray,
  boxrule=0.6pt,arc=1.5mm,title={#1},fonttitle=\bfseries}
\newtcolorbox{conceptcheck}[1][]{colback=white,colframe=blue,
  boxrule=0.6pt,arc=1.5mm,title={#1},fonttitle=\bfseries}

% Empty replacement box for missing blackboard drawings.
% Replace the contents of any instance with \includegraphics[width=\linewidth]{...}.
\newcommand{\boardplaceholder}[2][5.2cm]{%
  \begin{center}
  \fcolorbox{darkgray}{white}{%
    \begin{minipage}[c][#1][c]{0.88\linewidth}
      \centering\color{darkgray}
      \textbf{Blackboard illustration placeholder}\\[0.6em]
      #2\\[0.6em]
      \small Replace this box with the classroom image.
    \end{minipage}}
  \end{center}}

\lstset{basicstyle=\ttfamily\small,keywordstyle=\color{blue}\bfseries,
  commentstyle=\color{darkgray},frame=single,rulecolor=\color{darkgray},
  columns=fullflexible,showstringspaces=false}

\title{\vspace{-2em}\textbf{Lecture 1: Computation and Abstraction}}
\author{Dr. Changlai Du}
\date{}

\begin{document}
\maketitle

\begin{keyidea}[Learning goals]
After this lecture, you should be able to explain what an algorithm is, distinguish a function from an algorithm, describe the Turing-machine model and the universal Turing machine, state the Church--Turing thesis, compare algorithms through time and space complexity, and explain why hierarchical abstraction makes modern computing systems manageable.
\end{keyidea}

\tableofcontents

\section{Computation as a Mechanical Process}

\subsection{A motivating problem}

Consider the decision problem
\[
  \text{Is }35{,}500{,}001\text{ divisible by }113?
\]
Long division gives
\[
  35{,}500{,}001=113\times314{,}159+34.
\]
The remainder is $34$, so the answer is \emph{no}.

The familiar long-division procedure consists of a small collection of steps:
\begin{enumerate}
  \item divide or estimate the next quotient digit;
  \item multiply the divisor by that digit;
  \item subtract;
  \item bring down the next input digit;
  \item repeat until no input remains.
\end{enumerate}
Nothing in these steps requires a new creative insight once the procedure has been specified. A person or a machine can follow the same rules. This repeatable, rule-governed character is the first sense in which computation is \emph{mechanical}.

\boardplaceholder[6.0cm]{Long division of $35{,}500{,}001$ by $113$, showing quotient $314{,}159$ and remainder $34$.}

\subsection{Algorithm, input, and output}

\begin{definition}[Algorithm]
An \emph{algorithm} is a finite, unambiguous sequence of instructions for solving a class of problems. It specifies how to transform an input into an output.
\end{definition}

For the divisibility example, an input may consist of a dividend $a$ and a divisor $b$. The requested output determines the type of problem:

\begin{center}
\begin{tabularx}{0.92\linewidth}{>{\bfseries}p{0.23\linewidth}X p{0.27\linewidth}}
\toprule
Problem type & Required output & Example \\
\midrule
Decision & yes or no & Is $a$ divisible by $b$? \\
Search & a particular object or value & What is the remainder of $a/b$? \\
Counting & the number of valid objects & How many divisors does $a$ have? \\
Optimization & the best feasible object or value & Which valid solution has minimum cost? \\
\bottomrule
\end{tabularx}
\end{center}

The output $34$ solves the remainder-search problem. The output \emph{no} solves the original decision problem. Keeping the problem specification separate from the method prevents a common conceptual mistake.

\subsection{Correctness and efficiency}

The same function can have several algorithms. Repeatedly subtracting $113$ from $35{,}500{,}001$ will eventually produce the correct remainder, but it needs roughly $314{,}159$ subtractions. Long division processes the input one digit at a time and therefore requires a number of rounds proportional to the number of digits.

This example separates two questions:
\begin{itemize}
  \item \textbf{Computability:} Can any mechanical procedure solve the problem for every valid input?
  \item \textbf{Complexity:} How much time and memory does a particular procedure require as the input grows?
\end{itemize}

\section{From a Human Calculator to a Turing Machine}

When performing long division by hand, we use paper as memory, symbols for data, eyes to read, a hand and pencil to write, a finite collection of remembered instructions, and rules that determine the next action. Turing's model keeps only these essential components.

\subsection{The abstract machine}

A Turing machine has:
\begin{itemize}
  \item an unbounded one-dimensional tape divided into cells;
  \item a finite tape alphabet, including a blank symbol;
  \item a read/write head that scans one cell at a time;
  \item a finite set of control states, including a start state and one or more halting states;
  \item a transition function that determines the next action.
\end{itemize}

The transition function has the form
\[
  \delta(q,a)=(q',b,D),
\]
where $q$ is the current state, $a$ is the scanned symbol, $q'$ is the next state, $b$ is the symbol written to the tape, and $D\in\{L,R\}$ is the direction of movement.

\boardplaceholder[6.4cm]{Turing machine: an unbounded tape, a read/write head, and a finite-state controller. Include labels for read, write, move left/right, and machine state.}

\subsection{Why one-dimensional and unbounded?}

A sheet of paper is two-dimensional, but its cells can be enumerated and placed along a one-dimensional tape. This transformation simplifies the mechanism: the head needs only to move left or right. The precise enumeration is unimportant as long as it preserves all information.

The tape is \emph{unbounded} because the size of a valid input has no fixed upper limit. For any particular run, the machine uses only finitely many cells. If it needs more, the model permits more tape. This idealization separates computability from the accidental memory limit of a physical computer.

The controller has only \emph{finitely many} states. An infinitely complicated controller could simply encode an answer for every input, turning computation into a lookup table. Finite control forces the machine to obtain answers by applying reusable rules.

\boardplaceholder[4.8cm]{Classroom mapping from a two-dimensional grid of paper cells to a one-dimensional tape.}

\subsection{A Turing-machine view of long division}

Fix the divisor at $113$. A state can record the current remainder $r$, where $0\le r<113$. When the head reads the next decimal digit $d$, compute
\[
  t=10r+d,\qquad c=\left\lfloor\frac{t}{113}\right\rfloor,
  \qquad r'=t\bmod113.
\]
The machine writes the quotient digit $c$, moves right, and changes to the state representing $r'$. For the input $35500001$, the computation is:

\begin{center}
\small
\begin{tabular}{c@{\qquad}c@{\qquad}c}
\toprule
Digit read & Quotient prefix & Remainder state \\
\midrule
$3$ & $0$ & $3$ \\
$5$ & $00$ & $35$ \\
$5$ & $003$ & $16$ \\
$0$ & $0031$ & $47$ \\
$0$ & $00314$ & $18$ \\
$0$ & $003141$ & $67$ \\
$0$ & $0031415$ & $105$ \\
$1$ & $00314159$ & $34$ \\
\bottomrule
\end{tabular}
\end{center}

After the machine reaches the blank cell, it halts. The final non-leading-zero quotient is $314159$ and the remainder state is $34$.

\begin{keyidea}[What the model establishes]
The Turing machine is not a blueprint for a fast computer. It is a deliberately minimal mathematical model that makes the phrase ``mechanical procedure'' precise enough to study.
\end{keyidea}

\section{More Algorithms and More Machines}

\subsection{The Euclidean algorithm}

To find $\gcd(404,96)$, repeatedly replace the larger pair by the divisor and the remainder:
\begin{align*}
404 &= 4\cdot96+20,\\
 96 &= 4\cdot20+16,\\
 20 &= 1\cdot16+4,\\
 16 &= 4\cdot4+0.
\end{align*}
The last nonzero remainder is $4$, so $\gcd(404,96)=4$.

The Euclidean algorithm can be summarized as follows:
\begin{lstlisting}[language=Python]
def gcd(a, b):
    while b != 0:
        a, b = b, a % b
    return a
\end{lstlisting}

Like long division, this procedure uses a fixed collection of operations and a repeat rule. We can therefore design another Turing machine whose states implement it. Prime factorization offers a different algorithm for the same function, usually with a very different cost.

\boardplaceholder[5.2cm]{Euclidean algorithm for $404$ and $96$, with the successive pairs $(404,96)$, $(96,20)$, $(20,16)$, $(16,4)$.}

\subsection{Functions versus algorithms}

\begin{definition}[Function and algorithm]
A \emph{function} specifies what output corresponds to each input. An \emph{algorithm} specifies how to produce that output. Several algorithms may compute the same function.
\end{definition}

For example, $\gcd(a,b)$ is a function. The Euclidean algorithm and a factorization-based method are algorithms that compute it. Integer multiplication is also a function. The schoolbook multiplication algorithm for two $n$-digit integers uses $O(n^2)$ single-digit multiplications, while faster multiplication algorithms compute the same function with a lower asymptotic running time.

\section{The Universal Turing Machine}

Building a separate physical machine for every algorithm would be possible in principle but impractical. A universal Turing machine $U$ accepts both an encoded machine $\langle M\rangle$ and an input $x$:
\[
  U(\langle M\rangle,x)=M(x).
\]
It reads the description of $M$ and simulates $M$ step by step on $x$. Thus, one fixed machine can perform any computation described by another Turing machine.

\boardplaceholder[5.8cm]{Universal Turing machine with input $\langle M\rangle$ and $x$, showing $U$ simulating $M(x)$.}

The key idea is that a program and ordinary data can both be represented as symbols. A program can therefore be stored, copied, inspected, interpreted, compiled, or generated like other data. This \emph{program-as-data} principle underlies stored-program computers, operating systems, compilers, interpreters, program analyzers, and coding tools.

\section{Stored-Program Computers}

In a stored-program computer, memory contains both instructions and data. A CPU repeatedly performs a fetch--decode--execute cycle:
\begin{enumerate}
  \item fetch the instruction whose address is in the program counter;
  \item decode the instruction;
  \item execute it and update registers or memory;
  \item update the program counter.
\end{enumerate}
Conditional branches and jumps change the program counter, making selection and repetition possible.

The von Neumann architecture is a prominent organization of this idea: a processor communicates with a shared memory that stores both code and data. It is a practical architecture, while a Turing machine is a mathematical model. Both embody finite control operating on expandable memory.

\boardplaceholder[6.0cm]{Von Neumann architecture: CPU with control unit, ALU, registers and program counter; shared memory containing code and data; input/output connections.}

\section{The Church--Turing Thesis}

\begin{keyidea}[Church--Turing thesis]
Every function that is effectively computable by an algorithm can be computed by a Turing machine.
\end{keyidea}

Here, \emph{effectively computable} means computable by a mechanical, rule-following procedure. The thesis identifies this informal idea with the mathematically defined Turing-machine model.

It is a \emph{thesis}, not a theorem, because ``effectively computable'' is an informal concept rather than a separately formalized mathematical definition. No proof can connect a formal model to an informal notion in the ordinary theorem-proof sense. The thesis receives strong support from the fact that many independently proposed models, including lambda calculus, register machines, general-purpose programming languages, and idealized computers with unbounded memory, have the same computational power through mutual simulation.

A system is called \emph{Turing complete} when it can simulate a universal Turing machine, subject to the same idealization of unbounded memory and execution time. Finite control, memory that can grow without a fixed bound, and conditional looping provide the essential ingredients.

\subsection{What would challenge the thesis?}

A genuine counterexample would require an effective, physically realizable procedure that computes a function no Turing machine can compute. Proposals involving infinite precision or exotic space-time models challenge the physical assumptions. Ordinary quantum computers do not currently provide such a counterexample: they may change efficiency, but standard quantum computation remains within the class of Turing-computable functions.

\section{Limits of Computation}

The Church--Turing thesis characterizes the power of mechanical computation, but it does not say that every clearly stated problem has an algorithmic solution.

\subsection{Computable numbers}

A real number is computable if a finite program can generate its digits to any requested precision. There are only countably many finite programs but uncountably many real numbers. Therefore, most real numbers are not computable, even though familiar constants such as $\pi$ and $e$ are computable.

\subsection{The halting problem}

Suppose a universal procedure $H(P,x)$ could always determine whether program $P$ eventually halts on input $x$. Construct a program $D$ that uses $H$ on its own description:
\begin{lstlisting}[language=Python]
def D():
    if H(D, D):
        loop_forever()
    else:
        return
\end{lstlisting}
If $H$ predicts that $D$ halts, $D$ loops. If $H$ predicts that $D$ loops, $D$ halts. Either answer contradicts the prediction. Therefore, no such universal halting decider exists.

This result is a consequence of the same program-as-data principle that makes universal computation possible. A sufficiently general computer can receive descriptions of programs, including descriptions that refer to themselves.

\begin{keyidea}[The price of universality]
Universal computers can simulate arbitrary programs, but no universal program can perfectly predict every semantic property of arbitrary programs. Rice's theorem generalizes this limitation: every nontrivial semantic property of the function computed by an arbitrary program is undecidable.
\end{keyidea}

\section{Complexity: The Cost of Computation}

Computability asks whether some algorithm exists. Complexity asks how resources grow with input size. The most common resources are running time and memory.

\subsection{Input size and asymptotic growth}

For an integer, input size usually means the number of digits or bits, not the integer's numerical value. If a decimal integer has $n$ digits, a digit-by-digit algorithm naturally performs work as a function of $n$.

Big-$O$ notation describes an asymptotic upper bound. It suppresses constant factors, lower-order terms, exact processor speed, and behavior on small inputs. Common growth rates include
\[
 O(1),\quad O(\log n),\quad O(n),\quad O(n\log n),\quad
 O(n^2),\quad O(2^n).
\]
For sufficiently large $n$, the shape of the growth function matters more than small constant differences.

\subsection{Searching an attendance list}

Finding a name in an unsorted list by checking entries in order takes $O(n)$ time. If the list is sorted, binary search takes $O(\log n)$ comparisons: about $7$ comparisons for roughly $100$ entries and about $30$ for one billion entries. A hash table can provide $O(1)$ expected lookup time by using additional memory and preprocessing.

This example illustrates a recurring time--space tradeoff: a method may run faster because it stores a more helpful representation of the data.

\subsection{Fibonacci numbers}

The recurrence
\[
F(0)=0,\qquad F(1)=1,\qquad F(n)=F(n-1)+F(n-2)
\]
does not determine a unique algorithm.

\begin{center}
\begin{tabularx}{0.94\linewidth}{>{\bfseries}p{0.28\linewidth}p{0.20\linewidth}p{0.20\linewidth}X}
\toprule
Method & Time & Extra space & Main idea \\
\midrule
Naive recursion & $O(2^n)$ & $O(n)$ & Recomputes the same subproblems \\
Dynamic programming & $O(n)$ & $O(n)$ & Stores all previously computed values \\
Iterative & $O(n)$ & $O(1)$ & Keeps only the two latest values \\
\bottomrule
\end{tabularx}
\end{center}

\boardplaceholder[5.5cm]{Recursion tree for naive Fibonacci, highlighting repeated calls such as $F(n-2)$.}

\subsection{P and NP}

\begin{definition}[P and NP]
\textbf{P} is the class of decision problems solvable in polynomial time. \textbf{NP} is the class of decision problems for which a proposed yes-certificate can be verified in polynomial time.
\end{definition}

Every problem in P is in NP because a problem that can be solved efficiently can also have its answer verified efficiently. Whether every NP problem is also in P remains unknown:
\[
  \mathrm{P}\stackrel{?}{=}\mathrm{NP}.
\]

Subset Sum illustrates the distinction. Given integers and a target, the decision problem asks whether some subset sums to the target. A proposed subset is easy to check, but naive search examines as many as $2^n$ subsets. This does not prove that the problem lacks a polynomial-time algorithm; it explains why verification may appear much easier than discovery.

\section{Hierarchical Abstraction}

The Turing machine tells us what a mechanical computer can compute. Computer architecture explains how a physical machine executes instructions. Abstraction explains how people build usable systems on top of that machinery.

\begin{definition}[Abstraction and interface]
An \emph{abstraction} hides details that are unnecessary for the current task. An \emph{interface} specifies what a component makes available to its users without requiring them to know its internal implementation.
\end{definition}

When writing \texttt{2 + 3} in Python, a programmer does not need to reason about transistors. The operation passes through layers such as language semantics, interpreter or compiler, machine instructions, digital circuits, logic gates, transistors, and electrical behavior. Each layer exposes a simpler interface to the layer above it.

Likewise, clicking \emph{Send} in a messaging application may involve the application interface, operating-system services, network protocols, device drivers, hardware, and physical signals. No single application developer needs to rebuild all of these layers.

\boardplaceholder[7.0cm]{Hierarchical abstraction stack, for example: user $\rightarrow$ application $\rightarrow$ operating system $\rightarrow$ network or instruction-set layer $\rightarrow$ hardware $\rightarrow$ physical signals.}

\subsection{Leaky abstractions}

An abstraction leaks when lower-level details affect behavior visible at a higher level. Examples include:
\begin{itemize}
  \item a Python list appears able to grow, but physical memory remains finite;
  \item a network request resembles a function call, but latency and failure remain possible;
  \item a file appears to be a stable sequence of bytes, but storage devices can fail;
  \item two algorithms implement the same function, but their running times may differ dramatically.
\end{itemize}
Abstraction does not erase lower layers. It allows us to ignore them until performance, correctness, security, or failure forces us to cross an interface and inspect the implementation.

\section{Concept Checklist}

\begin{conceptcheck}[You should be able to explain each pair or relationship]
\begin{itemize}
  \item input, output, instruction, algorithm, and function;
  \item decision, search, counting, and optimization problems;
  \item model, simulation, mechanical procedure, and computability;
  \item tape, state, transition function, and halting state;
  \item Turing machine, universal Turing machine, and Turing completeness;
  \item stored-program computer, von Neumann architecture, and program counter;
  \item thesis versus theorem, and the Church--Turing thesis;
  \item computability versus decidability, including the halting problem;
  \item input size, Big-$O$, time complexity, and space complexity;
  \item recursion, dynamic programming, and iterative computation;
  \item P, NP, efficient solution, and efficient verification;
  \item abstraction, interface, implementation, and leaky abstraction.
\end{itemize}
\end{conceptcheck}

\section{Study Questions}

\begin{enumerate}
  \item Why does the long-division example qualify as a mechanical procedure?
  \item What changes when divisibility is posed as a decision problem rather than a remainder-search problem?
  \item Why does a Turing machine use an unbounded tape but only finitely many control states?
  \item Using the state-update rule $t=10r+d$, reproduce the remainder sequence for $35{,}500{,}001/113$.
  \item Explain why the Euclidean algorithm and prime factorization may compute the same GCD function while having different costs.
  \item What information must a universal Turing machine receive to simulate another machine?
  \item Why is the Church--Turing statement called a thesis rather than a theorem?
  \item How does program-as-data enable both stored-program computing and the halting-problem proof?
  \item Compare linear search, binary search, and hash lookup. What assumptions and tradeoffs distinguish them?
  \item Why do naive recursion and iteration have different costs for Fibonacci numbers even though both compute $F(n)$?
  \item State the P versus NP question without saying that NP means ``non-polynomial.''
  \item Give one example of an abstraction leak and identify the lower-level detail that becomes visible.
\end{enumerate}

\section*{Homework}
\addcontentsline{toc}{section}{Homework}
Watch the CS50 Week 0 video and complete all assigned parts:
\begin{center}
\url{https://cs50.harvard.edu/x/weeks/0/}
\end{center}

\section*{Sources Used for These Notes}
\addcontentsline{toc}{section}{Sources Used for These Notes}
These notes synthesize the slide deck \emph{01 Computation and Abstraction.pptx}, its embedded speaker notes, and the two classroom transcripts dated 8 September 2026. The wording has been edited for clarity, and transcription errors have been corrected using the mathematical context. Blackboard drawings that were described in speech but unavailable as images remain as labeled replacement boxes.

For historical context on the Turing machine, see A. M. Turing, ``On Computable Numbers, with an Application to the Entscheidungsproblem,'' \emph{Proceedings of the London Mathematical Society}, 1936. For a detailed discussion of the Church--Turing thesis, see \url{https://plato.stanford.edu/entries/church-turing/}.

\end{document}
